% !PRINTED 350
\weberpage{350}
%<<CONT

\[
\begin{aligned}
&\text{a)}\quad x^3 + e x - g = 0\\
&\text{b)}\quad x^3 - e x - g = 0, &&\frac{e^3}{27} < \frac{g^2}{4}\\
&\text{c)}\quad x^3 - e x - g = 0, &&\frac{e^3}{27} > \frac{g^2}{4}.
\end{aligned}
\]

We exclude the limiting cases in which one of the quantities $e$, $g$, $4e^3 - 27g^2$ vanishes.

To apply the \gesperrt{Cardan} formula (\S.\,35), we set

\weq{2}{R = \frac{g^2}{4} \pm \frac{e^3}{27},}

where the upper sign holds in case a), the lower in cases b) and c). In cases a) and b) only one root is real; in case c) (the Casus irreducibilis) all three roots are real. The \gesperrt{Cardan} formula gives

\weq{3}{x = \sqrt[3]{\frac{g}{2} + \sqrt{R}} + \sqrt[3]{\frac{g}{2} - \sqrt{R}}.}

We now proceed in the three cases as follows:

a) We introduce an angle $\vartheta$, which we can choose between $0^0$ and $90^0$, by the equation:

\weq{4}{\frac{g}{2} = \sqrt{\frac{e^3}{27}}\,\mathrm{cotg}\,\vartheta,}

hence

\[
\sqrt{R} = \sqrt{\frac{e^3}{27}}\,\frac{1}{\sin\vartheta},
\]

\[
\begin{aligned}
x &= \sqrt{\frac{e}{3}}\left(\sqrt[3]{\mathrm{cotg}\,\vartheta + \frac{1}{\sin\vartheta}} + \sqrt[3]{\mathrm{cotg}\,\vartheta - \frac{1}{\sin\vartheta}}\right)\\
&= \sqrt{\frac{e}{3}}\left(\sqrt[3]{\mathrm{cotg}\,\frac{\vartheta}{2}} - \sqrt[3]{\mathrm{tg}\,\frac{\vartheta}{2}}\right),
\end{aligned}
\]

and if we further determine an angle $\varphi$ from

\weq{5}{\mathrm{tg}\,\frac{\vartheta}{2} = \mathrm{tg}\,\varphi^3}

then we obtain

\weq{6}{x = 2\sqrt{\frac{e}{3}}\,\mathrm{cotg}\,2\varphi.}

The angles $\vartheta$, $\varphi$ and finally $x$ are found from the logarithmic-trigonometric tables by the formulas (4), (5), (6).
